\chapter{The Complete Pipeline}
\label{ch:complete-pipeline}

This chapter assembles the data sources (Chapter~\ref{ch:our-data}), feature layers (Chapters~\ref{ch:har-core}--\ref{ch:feature-composition}), models (Chapters~\ref{ch:lightgbm}--\ref{ch:rashomon}), and evaluation framework (Chapter~\ref{ch:evaluation}) into the end-to-end system and defines the implementation roadmap.

%% ──────────────────────────────────────────────
\section{End-to-End System Diagram}
\label{sec:system-diagram}
%% ──────────────────────────────────────────────

Figure~\ref{fig:full-pipeline} shows every component from raw data to scored forecasts.

\begin{figure}[htbp]
\centering
\begin{tikzpicture}[
  node distance=0.55cm and 0.35cm,
  >=Stealth,
  % Data sources: blue
  datasrc/.style={draw, rounded corners, minimum height=0.7cm, minimum width=2.1cm,
    align=center, font=\scriptsize, fill=blue!8, text width=2.0cm},
  % Computation: green
  compute/.style={draw, rounded corners, minimum height=0.7cm, minimum width=2.8cm,
    align=center, font=\scriptsize, fill=green!8},
  % Feature store: green
  store/.style={draw, rounded corners, minimum height=0.7cm, minimum width=4.5cm,
    align=center, font=\small, fill=green!8},
  % Models: orange
  model/.style={draw, rounded corners, minimum height=0.7cm, minimum width=2.6cm,
    align=center, font=\scriptsize, fill=orange!8},
  % Ensemble/output: orange
  blend/.style={draw, rounded corners, minimum height=0.7cm, minimum width=3.0cm,
    align=center, font=\small, fill=orange!8},
  % Evaluation: red
  evalbox/.style={draw, rounded corners, minimum height=0.7cm, minimum width=3.6cm,
    align=center, font=\small, fill=red!8},
  % Arrows
  arrow/.style={-{Stealth[length=2mm]}, thick},
  dasharrow/.style={-{Stealth[length=2mm]}, thick, dashed, gray},
]

% ── Row 1: Data Sources ──
\node[datasrc] (tick) {Tick-level\\$\RV$, RQ};
\node[datasrc, right=of tick] (ohlcv) {Daily\\OHLCV};
\node[datasrc, right=of ohlcv] (emini) {E-mini\\L2 Depth};
\node[datasrc, right=of emini] (iv) {SPX IV\\Surface};
\node[datasrc, right=of iv] (vix) {VIX Term\\Structure};
\node[datasrc, right=of vix] (cross) {Cross-\\Asset};

% ── Row 2: Feature Computation ──
\node[compute, below=0.8cm of $(tick)!0.5!(ohlcv)$] (feat01) {Layers 0--1\\HAR + Jumps};
\node[compute, below=0.8cm of emini] (feat3) {Layer 3\\Microstructure};
\node[compute, below=0.8cm of $(iv)!0.5!(vix)$] (feat2) {Layer 2\\Options};
\node[compute, below=0.8cm of cross] (feat4) {Layers 4--7\\Cross/Cal/Mem};

% Arrows: data -> computation
\draw[arrow] (tick.south) -- (feat01.north);
\draw[arrow] (ohlcv.south) -- (feat01.north);
\draw[arrow] (emini.south) -- (feat3.north);
\draw[arrow] (iv.south) -- (feat2.north);
\draw[arrow] (vix.south) -- (feat2.north);
\draw[arrow] (cross.south) -- (feat4.north);

% ── Row 3: Feature Store ──
\node[store, below=0.8cm of $(feat01)!0.5!(feat4)$] (fstore) {Feature Store ($\sim$80--120 features)};

% Arrows: computation -> store
\draw[arrow] (feat01.south) |- (fstore.north -| feat01.south);
\draw[arrow] (feat3.south) -- (fstore.north -| feat3.south);
\draw[arrow] (feat2.south) |- (fstore.north -| feat2.south);
\draw[arrow] (feat4.south) |- (fstore.north -| feat4.south);

% ── Row 4: Three model branches ──
\node[model, below left=0.9cm and 1.8cm of fstore] (lgbm) {LightGBM\\(Ch.~\ref{ch:lightgbm})};
\node[model, below=0.9cm of fstore] (lstm) {LSTM / TCN\\(Ch.~10)};
\node[model, below right=0.9cm and 1.8cm of fstore] (otree) {Optimal Trees\\(Ch.~\ref{ch:rashomon})};

% Arrows: store -> models
\draw[arrow] (fstore.south) -- ++(0,-0.3) -| (lgbm.north)
  node[pos=0.75, left, font=\tiny] {tabular};
\draw[arrow] (fstore.south) -- (lstm.north)
  node[midway, right, font=\tiny] {intraday seq.};
\draw[arrow] (fstore.south) -- ++(0,-0.3) -| (otree.north)
  node[pos=0.75, right, font=\tiny] {tabular};

% Also draw direct line from E-mini data to LSTM (raw sequences)
\draw[dasharrow] (emini.south east) -- ++(0.3,-0.3) |- ([yshift=0.2cm]lstm.north)
  node[pos=0.2, right, font=\tiny, gray] {raw bars};

% ── Row 5: Ensemble ──
\node[blend, below=0.9cm of lstm] (ensemble) {Ensemble Blend (Ch.~\ref{ch:ensemble})};

\draw[arrow] (lgbm.south) |- (ensemble.west);
\draw[arrow] (lstm.south) -- (ensemble.north);
\draw[arrow] (otree.south) |- (ensemble.east);

% ── Row 6: Forecast ──
\node[blend, below=0.7cm of ensemble] (forecast) {Forecast $\log \RV_{t+h}$};
\draw[arrow] (ensemble.south) -- (forecast.north);

% ── Row 7: Evaluation ──
\node[evalbox, below=0.7cm of forecast] (eval) {Evaluation: $\QLIKE$, DM, MCS (Ch.~\ref{ch:evaluation})};
\draw[arrow] (forecast.south) -- (eval.north);

\end{tikzpicture}
\caption{End-to-end pipeline from raw data to scored forecasts.
Blue: data sources. Green: feature computation and storage. Orange: models and blending. Red: evaluation.
Dashed arrow indicates the LSTM also receives raw intraday bar sequences directly.}
\label{fig:full-pipeline}
\end{figure}


%% ──────────────────────────────────────────────
\section{Implementation Order}
\label{sec:impl-order}
%% ──────────────────────────────────────────────

Each step produces a standalone, reportable result.
No step depends on a later step being complete.

\begin{table}[htbp]
\centering
\caption{Implementation roadmap. Each step is independently reportable.}
\label{tab:impl-order}
\small
\begin{tabular}{@{}clp{2.4cm}lp{3.4cm}@{}}
\toprule
\textbf{Step} & \textbf{What} & \textbf{Features} & \textbf{Model} & \textbf{Deliverable} \\
\midrule
1 & HARQ + SHAR baseline & Layers 0--1 (11 features) & OLS / Ridge & Walk-forward $\QLIKE$; baseline table \\
2 & Add options layer & Layer 2 ($\sim$20 total) & LightGBM & $\QLIKE$ lift vs.\ Step 1; $\SHAP$ summary \\
3 & Cross-asset + spillover & Layer 4 ($\sim$30 total) & LightGBM & Spillover contribution analysis \\
4 & E-mini microstructure & Layer 3 (separate) & LSTM / TCN & Standalone intraday forecast \\
5 & Full feature set & Layers 5--7 ($\sim$80--120) & Ensemble & Final $\QLIKE$; DM tests; MCS \\
6 & Rashomon analysis & Same as Step 5 & Optimal Trees & Variable importance stability report \\
\bottomrule
\end{tabular}
\end{table}

\begin{keyidea}[Steps 1--2 Are the Critical Path]
If Step~1 (HARQ/SHAR baseline) is weak, every subsequent $\QLIKE$ comparison is misleading \citep{HARdToBeat2024}.
Step~2 (adding options features via LightGBM) produces the first genuine ML-vs-baseline comparison.
Getting these two steps right matters more than reaching Step~5.
\end{keyidea}


%% ──────────────────────────────────────────────
\section{Re-training and Monitoring}
\label{sec:retrain}
%% ──────────────────────────────────────────────

Retrain all models weekly on a rolling 5-year window.
After each retrain, compute the Rashomon set and record the Jaccard similarity of the top-10 features compared to the previous window; a drop below 0.6 signals feature-importance drift and warrants investigation.
Monitor out-of-sample $\QLIKE$ on a trailing 20-day basis; alert if it degrades more than 10\% relative to the trailing 60-day average.
Log every retrain's feature ranking and $\QLIKE$ for post-hoc regime analysis.


%% ──────────────────────────────────────────────
\section{Lookahead Bias Checklist}
\label{sec:lookahead}
%% ──────────────────────────────────────────────

Table~\ref{tab:lookahead} lists the four primary sources of lookahead bias in volatility forecasting systems.

\begin{table}[htbp]
\centering
\caption{Lookahead bias sources and prevention rules.}
\label{tab:lookahead}
\small
\begin{tabular}{@{}lp{4.0cm}p{5.2cm}@{}}
\toprule
\textbf{Source} & \textbf{Pitfall} & \textbf{Rule} \\
\midrule
Realized measures & Intraday returns from the target day leak into features & Features for $\RV_{t+1}$ use only information $\leq t$ \\
Microstructure    & Full-day LOB features include the close & Truncate intraday sequences at $t-\epsilon$; align timestamps strictly \\
Options surface   & Intraday surface changes reflect information about the target day & Use end-of-day surface from day $t$ for the day-$t{+}1$ prediction \\
Cross-asset       & Mixed frequencies across asset classes & Synchronize all inputs to the same end-of-day cutoff \\
\bottomrule
\end{tabular}
\end{table}

\begin{warning}[Lookahead Bias]
The single most common error in financial ML research.
Every feature must be computable strictly before the forecast target period begins.
When in doubt, add a one-day lag.
A lookahead-contaminated model will show excellent in-sample $\QLIKE$ that vanishes out of sample, wasting weeks of development time before the bug is identified.
\end{warning}
